\documentclass[10pt,a4paper]{amsart}
\usepackage{amsmath,amssymb,amsthm,graphicx}
\usepackage{verbatim}
\newenvironment{reference}{\comment}{\endcomment}
%\newenvironment{reference}{}{}
\newenvironment{slogan}{\comment}{\endcomment}
\newenvironment{history}{\comment}{\endcomment}

% For commutative diagrams we use Xy-pic
\usepackage[all]{xy}

% We use 2cell for 2-commutative diagrams.
\xyoption{2cell}
% Theorem environments.
%
\theoremstyle{plain}
\newtheorem{theorem}[subsection]{Theorem}
\newtheorem{proposition}[subsection]{Proposition}
\newtheorem{lemma}[subsection]{Lemma}

\theoremstyle{definition}
\newtheorem{definition}[subsection]{Definition}
\newtheorem{example}[subsection]{Example}
\newtheorem{exercise}[subsection]{Exercise}
\newtheorem{situation}[subsection]{Situation}

\theoremstyle{remark}
\newtheorem{remark}[subsection]{Remark}
\newtheorem{remarks}[subsection]{Remarks}

\numberwithin{equation}{subsection}

% Macros
%
\def\lim{\mathop{\mathrm{lim}}\nolimits}
\def\colim{\mathop{\mathrm{colim}}\nolimits}
\def\Spec{\mathop{\mathrm{Spec}}}
\def\Hom{\mathop{\mathrm{Hom}}\nolimits}
\def\Ext{\mathop{\mathrm{Ext}}\nolimits}
\def\SheafHom{\mathop{\mathcal{H}\!\mathit{om}}\nolimits}
\def\SheafExt{\mathop{\mathcal{E}\!\mathit{xt}}\nolimits}
\def\Sch{\mathit{Sch}}
\def\Mor{\mathop{\mathrm{Mor}}\nolimits}
\def\Ob{\mathop{\mathrm{Ob}}\nolimits}
\def\Sh{\mathop{\mathit{Sh}}\nolimits}
\def\NL{\mathop{N\!L}\nolimits}
\def\CH{\mathop{\mathrm{CH}}\nolimits}
\def\proetale{{pro\text{-}\acute{e}tale}}
\def\etale{{\acute{e}tale}}
\def\QCoh{\mathit{QCoh}}
\def\Ker{\mathop{\mathrm{Ker}}}
\def\Im{\mathop{\mathrm{Im}}}
\def\Coker{\mathop{\mathrm{Coker}}}
\def\Coim{\mathop{\mathrm{Coim}}}

% Boxtimes
%
\DeclareMathSymbol{\boxtimes}{\mathbin}{AMSa}{"02}

%
% Macros for moduli stacks/spaces
%
\def\QCohstack{\mathcal{QC}\!\mathit{oh}}
\def\Cohstack{\mathcal{C}\!\mathit{oh}}
\def\Spacesstack{\mathcal{S}\!\mathit{paces}}
\def\Quotfunctor{\mathrm{Quot}}
\def\Hilbfunctor{\mathrm{Hilb}}
\def\Curvesstack{\mathcal{C}\!\mathit{urves}}
\def\Polarizedstack{\mathcal{P}\!\mathit{olarized}}
\def\Complexesstack{\mathcal{C}\!\mathit{omplexes}}
% \Pic is the operator that assigns to X its picard group, usage \Pic(X)
% \Picardstack_{X/B} denotes the Picard stack of X over B
% \Picardfunctor_{X/B} denotes the Picard functor of X over B
\def\Pic{\mathop{\mathrm{Pic}}\nolimits}
\def\Picardstack{\mathcal{P}\!\mathit{ic}}
\def\Picardfunctor{\mathrm{Pic}}
\def\Deformationcategory{\mathcal{D}\!\mathit{ef}}
\usepackage{cleveref}
\crefname{originalreference}{original reference}{original references}
\Crefname{originalreference}{Original reference}{Original references}
\makeatletter
\newlabel{morphisms-lemma-flat-base-change-scheme-theoretic-image}{{081I}{source}{}{sourcepointer-1}{}}
\newlabel{morphisms-lemma-flat-base-change-scheme-theoretic-image@cref}{{[originalreference][0][]081I}{[1][0][]source}{}{}{}}
\makeatother
\crefname{originalreference}{original reference}{original references}
\Crefname{originalreference}{Original reference}{Original references}
\makeatletter
\newlabel{properties-lemma-complement-closed-point-Jacobson}{{02IM}{source}{}{sourcepointer-2}{}}
\newlabel{properties-lemma-complement-closed-point-Jacobson@cref}{{[originalreference][0][]02IM}{[1][0][]source}{}{}{}}
\makeatother
\title{Dimension is constant on nonempty opens of irreducible Jacobson catenary schemes}
\author{The Stacks Project Authors}
\date{stacks-project:stacks-geometry.tex:local-snapshot-20261012; GNU FDL 1.2 or later}
\begin{document}
\maketitle
\noindent\textit{Editorial note.} All delivered mathematical text is the author's own, copied verbatim from the pinned Stacks Project chapter file stacks-geometry.tex; no statement or proof step was rewritten, reordered or removed.\par\medskip
\noindent\textit{Editorial note.} Span recomputation by label pairing (LANE\_BRIEF 10.2): in this chapter the author's proof environment follows the statement environment as a standalone block; the spec covers the complete environment pair (statement span, proof span) as recorded in the pool candidate.\par\medskip
\noindent\textit{Editorial note.} Occupation scan: combined claim\_key + statement\_anchor + generic-line scan over all registration files (active and withdrawn) found 0 hits for stacks-geometry.tex (ge-occupied.json, 2026-10-14); 207 route=geometry / chapter-list field hits in focused-seeds are alias noise (route name and alphabetical chapter enumerations), excluded per the ladder alias-blind-spot caveat.\par\medskip
\ifdefined\hypertarget\hypertarget{sourcepointer-1}{}\fi
\noindent\textit{Source reference.} Original reference 081I: sources/stacks-chow/tags.txt, source lines 10218--10218. Cross-chapter standard prerequisite cited by the author inside the counted span; it is used as a pointer and its own proof is not part of this item.\par\medskip
\ifdefined\hypertarget\hypertarget{sourcepointer-2}{}\fi
\noindent\textit{Source reference.} Original reference 02IM: sources/stacks-chow/tags.txt, source lines 3305--3305. Cross-chapter standard prerequisite cited by the author inside the counted span; it is used as a pointer and its own proof is not part of this item.\par\medskip
\begin{lemma}
\label{lemma-constancy-of-dimension}
If $X$ is an irreducible, Jacobson, catenary, and locally Noetherian
scheme of finite dimension,
then $\dim U = \dim X$ for every
non-empty open subset $U$ of $X$.
Equivalently, $\dim_x X$ is a constant function on $X$.
\end{lemma}

\begin{proof}
The equivalence of the two claims follows directly from the
definitions.   Suppose, then, that $U\subset X$ is a non-empty open
subset.
Certainly $\dim U \leq \dim X$, and we have to show
that $\dim U \geq \dim X.$
Write $d = \dim X$, and choose a maximal strictly
decreasing chain of irreducible closed subsets
of $X$, say
$$
X = Z_0 \supset Z_1 \supset \ldots \supset Z_d.
$$
Since $X$ is Jacobson, the minimal irreducible closed
subset $Z_d$ is equal to $\{x\}$ for some closed
point $x$.

\medskip\noindent
If $x \in U,$ then
$$
U = U \cap Z_0  \supset U\cap Z_1 \supset \ldots \supset
U\cap Z_d = \{x\}
$$
is a strictly decreasing chain of irreducible closed
subsets of $U$, and so we conclude that $\dim U \geq d$,
as required.  Thus we may suppose that $x \not\in U.$

\medskip\noindent
Consider the flat morphism $\Spec \mathcal{O}_{X,x} \to X$.
The non-empty (and hence dense) open subset $U$ of $X$
pulls back to an open subset $V \subset \Spec \mathcal{O}_{X,x}$.
Replacing $U$ by a non-empty quasi-compact, and hence
Noetherian, open subset, we may assume that the inclusion
$U \to X$ is a quasi-compact morphism.  Since the
formation of scheme-theoretic images of quasi-compact
morphisms commutes with flat base-change
Morphisms, Lemma
\ref{morphisms-lemma-flat-base-change-scheme-theoretic-image}
we see that $V$ is dense in $\Spec \mathcal{O}_{X,x}$,
and so in particular non-empty,
and of course $x \not\in V.$  (Here we use $x$ also to denote
the closed point of $\Spec \mathcal{O}_{X,x}$, since its image
is equal to the given point $x \in X$.)
Now $\Spec \mathcal{O}_{X,x} \setminus \{x\}$ is Jacobson
Properties, Lemma
\ref{properties-lemma-complement-closed-point-Jacobson}
and hence $V$ contains a closed point $z$
of $\Spec \mathcal{O}_{X,x} \setminus \{x\}$.  The closure
in $X$ of the image of $z$ is then an irreducible
closed subset $Z$ of $X$ containing $x$, whose intersection
with $U$ is non-empty, and
for which there is no irreducible closed
subset properly contained in $Z$
and properly containing $\{x\}$
(because pull-back to $\Spec \mathcal{O}_{X,x}$ induces
a bijection between irreducible closed subsets of $X$
containing $x$ and irreducible closed subsets of $\Spec
\mathcal{O}_{X,x}$).
Since $U \cap Z$ is a non-empty closed subset of $U$,
it contains a point $u$ that is closed in $X$ (since
$X$ is Jacobson), and since $U\cap Z$
is a non-empty (and hence dense) open subset of the irreducible set $Z$
(which contains a point not lying in $U$, namely $x$),
the inclusion $\{u\} \subset U\cap Z$ is proper.

\medskip\noindent
As $X$ is catenary, the chain
$$
X = Z_0 \supset Z \supset \{x\} = Z_d
$$
can be refined to a chain of length $d+1$, which must then
be of the form
$$
X = Z_0 \supset W_1 \supset \ldots \supset W_{d-1} = Z \supset \{x\} = Z_d.
$$
Since $U\cap Z$ is non-empty, we then find that
$$
U = U \cap Z_0 \supset U \cap W_1\supset \ldots \supset U\cap W_{d-1}
= U\cap Z \supset \{u\}
$$
is a strictly decreasing chain of irreducible closed subsets
of $U$ of length $d+1$, showing that $\dim U \geq d$,
as required.
\end{proof}

\end{document}
